\begin{figure}[htbp]
\centering
\setlength{\fboxsep}{7pt}
\begin{tabular}{c c c}
\fbox{\begin{minipage}{0.24\linewidth}
\centering
\textbf{recipes}\\
\small id, name, category, difficulty, servings, instructions
\end{minipage}}
&
\fbox{\begin{minipage}{0.26\linewidth}
\centering
\textbf{recipe\_ingredients}\\
\small recipe\_id, ingredient\_id, quantity\_g
\end{minipage}}
&
\fbox{\begin{minipage}{0.28\linewidth}
\centering
\textbf{ingredients}\\
\small id, name, calories, protein, fat, carbs, allergen notes
\end{minipage}}
\\[0.3cm]
\multicolumn{3}{c}{\small many recipes can use many ingredients; the join table stores the quantity used}
\end{tabular}
\caption{Relational model showing why recipe nutrition can be derived from ingredient quantities.}
\label{fig:data-model}
\end{figure}
